\section{多 Agent 编排：从结构自由度到工程纪律}

\subsection{编排维度}
课程中对多 Agent 编排给出了清晰的设计维度：静态/动态、集中式/去中心化、上下文共享/隔离、自动化/人工介入。这些维度决定系统性能上限与稳定性下限。

\begin{figure}[H]
\centering
\includegraphics[width=0.93\textwidth]{frame-04-orchestration.jpg}
\caption{Multi-Agent Orchestration 设计维度}
\end{figure}
\noindent\footnotesize 画面证据：字幕约在 00:06:02--00:07:40，讲者进入复杂任务求解与轨迹组织问题。\normalsize

\begin{importantbox}{多 Agent 不是“多窗口并行”，而是“任务接口设计”}
如果没有角色边界、上下文边界与验收边界，增加 Agent 数量通常只会放大冲突与重复劳动。  
真正有效的并行来自清晰的输入输出契约与中间产物格式。
\end{importantbox}

\subsection{设计模式：Conversation / Tool Use / Planning / Memory}
\begin{figure}[H]
\centering
\includegraphics[width=0.93\textwidth]{frame-05-patterns.jpg}
\caption{Agentic Design Patterns：对话、工具、规划与记忆协同}
\end{figure}
\noindent\footnotesize 画面证据：字幕约在 00:07:50--00:09:10，讲者连续讨论 reflection、plan、fact-check 等模式。\normalsize

\begin{knowledgebox}{四类模式如何形成最小可用系统}
\begin{itemize}
    \item Conversation：承载角色协作与上下文传递；
    \item Tool Use：把语言决策映射到可执行动作；
    \item Planning：将长任务拆成可验证子目标；
    \item Memory：跨轮保留与压缩关键状态。
\end{itemize}
\end{knowledgebox}

\begin{warningbox}{只做“模式堆叠”会导致系统复杂度失控}
模式并非越多越好。每增加一种机制，都应回答两个问题：  
它降低了哪类失败率？它引入了哪类新故障模式？
\end{warningbox}

\subsection{本章小结}
多 Agent 工程的核心不是“创意编排”，而是“约束设计”。只有接口清晰、状态可检验、回滚可执行，复杂编排才可能稳定落地。


